\chapter{Introduction}

\section{The Past, Present, and Future of Common Lisp\index{Common Lisp}}

\subsection{Lisp Yesterday}

John McCarthy discovered the basic principles of Lisp in 1958, when he was processing
complex mathematical lists at MIT. Common Lisp (CL) is a high-level computer language,
whose syntax follows a simple list-like structure. The term ``Lisp'' itself originally stood for
``LISt Processing.'' When developing, testing, and running a CL program, at the core is a
modern-day version of the original List Processor which processes (compiles, evaluates, etc.)
the elements of your program. These elements, at the source code level, are represented
as lists. A list, in this context, is just a sequence of items, much like a familiar shopping
list or checklist. Originally the list was pretty much the only data structure supported by
Lisp, but modern-day Common Lisp supports a wide range of flexible and efficient data
structures.

A typical Common Lisp development and runtime environment behaves like a complete
operating system, with multiple threads of execution and the ability to load new code
and redefine objects (functions, etc.) dynamically, i.e. without stopping and restarting the
``machine.''

This ``operating system'' characteristic also makes CL significantly more flexible than
other popular programming languages. Whereas some other languages, such as shell scripts
or Perl CGI scripts\index{CGI scripts}, need to pipe data around, in CL all data manipulations can run
interactively within one single process with a shared memory space.

\subsection{Lisp Today}

The most popular form of Lisp used today, ``ANSI Common Lisp\index{ANSI Common Lisp},'' was initially designed by
Guy Steele in \emph{Common Lisp, the Language, 2nd Edition} (``CLtL2'') --- (see the Bibliography in Appendix B). This version of Lisp became the accepted industry standard. In 1995,
the American National Standards Institute\index{American National Standards Institute} recognized a slightly updated version as ANSI
Common Lisp, the first object-oriented language to receive certification, and ANSI CL remains the only language that meets all of the criteria set forth by the Object Management
Group (OMG) for a complete object-oriented language.

Paul Graham outlines ANSI CL in his 1996 book \emph{ANSI Common Lisp}, also listed in
Appendix B.

Official language standardization is important, because it protects developers from becoming saddled with legacy applications as new versions of a language are implemented.
For example, this lack of standardization has been a continuing problem for Perl developers.
Since each implementation of Perl defines the behavior of the language, it is not uncommon
to have applications that require outdated versions of Perl, making it nearly impossible to
use in a mission-critical commercial environment.

\subsection{Lisp Tomorrow}

Many developers once hoped that the software development process of the future would be
more automated through Computer-aided Software Engineering\index{Computer-aided Software Engineering} (CASE) tools. Such tools
claim to enable programmers and non-programmers to diagram their applications visually
and automatically generate code. While useful to a certain extent, traditional CASE tools
cannot cover domain-specific details and support all possible kinds of customizations --- so
developers inevitably need to hand-code the rest of their applications, breaking the link
between the CASE model and the code. CASE systems fall short of being a true ``problem-solving tool.''

The recursive nature of CL, and its natural ability to build applications in layers and
build upon itself --- in essence, code which writes code which writes code..., makes CL a
much better software engineering solution. CL programs can generate other CL programs
(usually at compile-time), allowing the developer to define all parts of the application in
one unified high-level language. Using macros and functions, naturally supported by the CL
syntax, the system can convert applications written in the high-level language automatically
into ``final'' code. Thus, CL is both a programming language and a powerful CASE tool,
and there is no need to break the connection.

Several other features of CL also save significant development time and manpower:

\begin{itemize}
\item \textbf{Automatic Memory Management\index{Automatic Memory Management}/Garbage Collection:} inherent in CL's basic architecture.

\item \textbf{Dynamic Typing\index{Dynamic Typing}:} In CL, values have types, but variables do not. This means that
you don't have to declare variables, or placeholders, ahead of time to be of a particular
type; you can simply create them or modify them on the fly\footnote{While you don't have to declare variable types, you may declare them. Doing so can help CL's compiler
to optimize your program further.}.

\item \textbf{Dynamic Redefinition\index{Dynamic Redefinition}:} You can add new operations and functionality to the running
CL process without the need for downtime. In fact, a program can be running in one
thread, while parts of the code are redefined in another thread. Moreover, as soon as
the redefinitions are finished, the application will use the new code and is effectively
modified while running. This feature is especially useful for server applications that
cannot afford downtimes -- you can patch and repair them while they are running.

\item \textbf{Portability:} The ``abstract machine'' quality of CL makes programs especially portable.
\end{itemize}

\begin{itemize}
\item \textbf{Standard Features:} CL contains many built-in features and data structures that are
non-standard in other languages. Features such as full-blown symbol tables and package systems, hash tables\index{hash tables}, list structures, and a comprehensive object system are several ``automatic'' features which would require significant developer time to recreate
in other languages.
\end{itemize}

\section{Convergence of Hardware and Software}

The most common criticism of Lisp usually made by programmers unfamiliar with the
language is that Lisp applications are too big and too slow. While this may have been true
20 years ago, it is absolutely not the case today. Today's inexpensive commodity computers
have more memory than some of the most powerful machines available just five years ago,
and easily meet the computing needs of a typical CL application. Hardware is no longer
the critical factor. Programmers and development time are!

Further, CL programming teams tend to be small, because the inherent features in CL
empower developers to do more. Without excessive effort, a single CL developer can create
and deploy a sophisticated, powerful program in a much shorter period of time than if
another language were used.

\section{The CL Model of Computation}

Programming in CL is distinguished from programming in other languages due to its unique
syntax and development mode. CL allows developers to make changes and test them immediately, in an incremental manner. In other languages, a developer must follow the
``compile-link-run-test'' cycle. These extra steps add minutes or hours to each cycle of development, and break the programmer's thought process. Multiply these figures by days,
weeks, and months --- and the potential savings are phenomenal.
